\chapter{Exchangeable Array Processes}

%%%%%%%%%%%%%%%%%%%%%%%%%%%%%%%%%%%%%%%%%%%%%%%%%%%%%%%%%%%%%%%%%%
% Introduction
%%%%%%%%%%%%%%%%%%%%%%%%%%%%%%%%%%%%%%%%%%%%%%%%%%%%%%%%%%%%%%%%%%
\section{Introduction}

The goal of this chapter is to develop a general framework for constructing and studying consistent
families of finite combinatorial structures, extending the program proposed in~\cite{aldous2010_more_uses_of_exchangeability}
toward a general theory of combinatorial limit objects. We aim to unify three ideas: Doob-Martin boundaries, graphons and hypergraphons,
and the classical theory of exchangeable arrays. By encoding the finite states of a growing combinatorial process as
restrictions of an exchangeable array, we simultaneously obtain a natural Markov chain and associated Martin boundary problem,
a sampling representation of its limiting objects, and a natural notion of finite-substructure embedding. We briefly discuss these three fields.


Exchangeability arises naturally when modeling random systems composed of many interchangeable parts.
The subject originates with de Finetti’s representation theorem for exchangeable binary sequences and was
subsequently extended to a wide variety of random combinatorial structures, including exchangeable
partitions and higher-dimensional arrays~\cite{kingman1978representation,austin2013_exchangeable_random_arrays,kallenberg2005_probabilistic_symmetries}.
A common feature of these theories is a representation theorem that generates an exchangeable
object by sampling independent latent random variables and applying an appropriate measurable function.
For the present chapter, the central result is the Aldous-Hoover-Kallenberg representation theorem for
exchangeable arrays, stated below as Theorem~\ref{thm:representation_theorem_of_exchangeable_arrays}.
This theorem places many random combinatorial structures within a common sampling framework and
reduces the construction of a corresponding limit theory to the problem of encoding their
finite substructures as entries of an exchangeable array.


Martin boundary theory, whose foundations were set in~\cite{doob1959_discrete_potential_theory_boundaries},
provides another perspective on this goal. Given a transient Markov chain $\Xn$ on some
countable state space $E$, Martin boundary theory aims to complete $E$ into a compact metric space
by adding ``points at $\infty$'' to which the Markov chain is converging. These boundary points
provide our first natural candidate for the desired limit objects, and this theory
provides a mechanism for tilting the underlying chain in a way that forces it to converge to
a given boundary point, analogous to how one can generate a sequence of graphs converging to a given
graphon via sampling. A short review of this topic and further references can be found
in Appendix~\ref{appendix:doob_martin_background_appendix}.

More recently, the theory of \textit{graphons} provided
a formalization for the limits of sequences of dense graphs. First developed in~\cite{lovasz2004_limits_dense_graph_sequences, borgs2007_convergent_dense_graph_sequences},
this theory topologizes the space of graphs by looking at a graph's finite substructures. In particular, two
graphs $G,H$ are said to be close if for each ``test'' graph $F$, the weak embedding density $emb^W(F,G)$ is close to
that of $H$, where
%
    \[emb^W(F,G) :=
    \frac{|\{\text{injective graph homomorphisms }\phi:F\to G\}|}
    {(|V(G)|)_{|V(F)|}}\]
%
is the fraction of injections $\phi:V(F) \to V(G)$ that preserve edges, and
$(N)_k:=N(N-1)\cdots(N-k+1)$. Since many other
combinatorial objects have an analogous notion of embedding,
this framework readily extends to other settings. A large body of work quickly
followed, extending these results to hypergraphs~\cite{zhao2015_hypergraph_limits_regularity_approach, elek2012_measure_theoretic_dense_hypergraphs},
 sparse graphs~\cite{borgs2019_sampling_sparse_exchangeable_graphs}, permutations~\cite{HoppenEtAl2013PermutationLimits},
and beyond.

Partial progress has already been made in connecting these strands of work.
The relationship between graphons and exchangeable arrays is well known: graphon
sampling produces a dissociated binary exchangeable array, the underlying adjacency matrix
of the graph sequence, while the corresponding
exchangeable array representation recovers the graphon up to the usual measure-preserving
equivalence~\cite{austin2013_exchangeable_random_arrays,tao2007correspondence}.
Moreover, Gr\"ubel and collaborators have shown that the usual graphon and permuton topologies agree
with the Doob-Martin topologies of corresponding random growth
chains~\cite{grubel2015_persisting_randomness_graphs_search_trees,grubel2023_ranks_copulas_permutons}.

These identifications, however, have generally been established separately for each
class of combinatorial structures. One first identifies the relevant finite substructures,
constructs an appropriate growth chain which preserves these structures, and then proves that its Doob-Martin topology agrees
with an independently defined limit topology. We aim to make this construction systematic. Once the
relationships among finite collections of components are encoded by an exchangeable array,
finite restrictions generate the labeled growth process, quotienting by relabeling generates the
corresponding unlabeled chain, and finite array patterns provide a natural notion of substructure density.\\


\noindent \textbf{Outline.} In Section~\ref{sec:exchangeable_array_processes}, we start by
defining two simple Markov processes associated to a given exchangeable array: a \textit{labeled}
chain generated by successively revealing the ``top left corner'' of the given array, and an \textit{unlabeled}
chain by quotienting the former up to relabeling the indices of the array. Furthermore, we discuss
how several existing examples of exchangeable combinatorial growth processes can be embedded into this framework.

In Section~\ref{sec:exchangeable_array_boundary}, we characterize the Doob-Martin boundary of the
labeled exchangeable array process and the minimal boundary of the unlabeled process. In the former, we show the boundary is the set of
deterministic arrays on which the exchangeable array is supported. In the latter, we show the minimal boundary
is characterized by the set of ergodic arrays. Equivalently, these arrays admit Aldous--Hoover--Kallenberg
representation functions that do not depend on the first coordinate, providing a natural analog of graphons and hypergraphons.

In Section~\ref{sec:embedding_density_representation_functions}, we define a few natural notions
of embedding a smaller (unlabeled) array into a larger one. As in the case of graphons, these embeddings provide
a way of topologizing the set of arrays: we say a sequence $y_n$ of arrays converges if the embedding density
$emb(x, y_n)$ converges for every finite array $x$. We show that these different notions of embedding
generate the same topology, and that this topology agrees with the Doob-Martin topology. We define a way
to extend these embedding densities to the limiting representation functions, and a way
of embedding a finite array into this space of representation functions in a way that
asymptotically preserves these densities.

Lastly, in Appendix~\ref{appendix:cut_metric_partial_progress}, we discuss some partial progress
on extending the classical cut metric from the theory of graphons to binary exchangeable array processes.
We introduce a strong cut norm for which an analog of the classical counting lemma holds, and exhibit a
random sequence that left-converges in probability while its canonical step-function representatives fail to be
Cauchy in probability in this norm. Furthermore,
we discuss how one might possibly extend a suitable metric from binary arrays to arrays with more general state spaces.

For clarity of exposition, most proofs are postponed until Appendix~\ref{appendix:proof_appendix}.





%%%%%%%%%%%%%%%%%%%%%%%%%%%%%%%%%%%%%%%%%%%%%%%%%%%%%%%%%%%%%%%%%%
% Exchangeable Array Processes
%%%%%%%%%%%%%%%%%%%%%%%%%%%%%%%%%%%%%%%%%%%%%%%%%%%%%%%%%%%%%%%%%%
\section{Exchangeable Array Processes}
\label{sec:exchangeable_array_processes}

Given a set $A$, let $\binom{A}{d}$ denote the set of all subsets of size $d$ of $A$, and let $[n] := \{1,2,\ldots,n\}$.
 A \textit{$d$-dimensional (jointly) exchangeable array} is a collection of random variables
 $(X_e)_{e \in \binom{\N}{d}}$ indexed by $d$-subsets of $\N$ whose distribution is invariant under
 simultaneous permutation of the coordinates. Namely, letting $\operatorname{Sym}(\N)$ be the set of finite
 permutations of $\N$ then
%
\[(X_{\sigma(e)})_{e \in \binom{\N}{d}} \overset{d}{=} (X_e)_{e \in \binom{\N}{d}}, \quad \forall \sigma \in \operatorname{Sym}(\N),\]
%
\noindent where $\sigma(e) = \{\sigma(i_1),\ldots,\sigma(i_d)\}$ for $e = \{i_1,\ldots,i_d\}$.

If $(X_e)_{e \in \binom{\N}{d}}$ is an exchangeable array there are two natural Markov processes
associated to it: a ``labeled'' chain $L_n$ and an ``unlabeled'' one $M_n$. Informally, the labeled
chain $L_n$ is produced by just revealing the entries with maximal coordinate $n$ at step $n$. Namely,
the finite labeled chain $L_n$ is just the ``top left corner'' of the array: the restriction
 of the array to entries corresponding to $\binom{[n]}{d}$. In this framework, the full
  array $(X_e)$ plays the role of the projective limit of the labeled chain which often
   appears in Doob--Martin theory.

The unlabeled chain $M_n$ is produced from the labeled one by ``forgetting the order'' of the rows,
columns, et cetera and just keeping track of the set of values appearing in each row, column, and
so on. More formally:

\begin{definition}[Exchangeable Array Processes]\label{def:exchangeable_array_process}
Given an exchangeable array $(X_e)_{e \in \binom{\N}{d}}$ with entries in some Borel
space $S$, define a graded Markov chain $L_n$ whose state space is
$E := \bigsqcup_{n \in \N} E_n$ for $E_n \subseteq S^{\binom{[n]}{d}}$ by
%
    \[(L_n)_e := X_e\quad \forall e \in \binom{[n]}{d}.\]
%
\noindent Define an equivalence relation $\sim$ on $E_n$ by saying
$A \sim B$ if there exists some $\sigma \in \operatorname{Sym}([n])$ so that $A_e = B_{\sigma(e)}$
for all $e \in \binom{[n]}{d}$, where $\sigma(e) = \{\sigma(i_1),\sigma(i_2),\ldots,\sigma(i_d)\}$
just permutes all elements of the subset $e$ by the same permutation. With this, define
the unlabeled chain
%
\[M_n := [L_n]_\sim.\]
%
\end{definition}
%
\noindent In this chapter, we will almost exclusively think of $S$ as countable,
and often finite. In order to apply Doob-Martin theory, we require each state in $E_n$
to be reachable from the empty array: $\P(L_n = x) > 0$ for all $x \in E_n$. In particular,
it suffices to choose $E_n = \mathrm{supp}(L_n)$.

Note that $L_n$ is trivially a Markov process in the sense
that any sequence of random variables becomes a Markov process if the state
space is just the entire history of the process. In general, projections of
Markov chains are not Markovian, but in this case, exchangeability ensures
that the Markov property is preserved.

% lemmma
\begin{lemma}[Markov Property of the Unlabeled Chain]
    When the state space $S$ is countable, the process $M_n$ is a Markov chain.

    \label{lem:index_equivalence}
\end{lemma}
%

\noindent The above lemma is true for any initial distribution, but we
are almost exclusively interested in the case where $M_0$ is the empty array.
The main benefit of this framework is that exchangeable arrays share a
strong structural property: they have
a representation through sampling from some hypercube.


\begin{theorem}[Representation Theorem]
\noindent An exchangeable array on $d$-subsets of $\N$ has a representation of the form

\[\left(X_e\right)_{e \in \binom{\N}{d}} \overset{d}{=}
\left(f\left(U_\emptyset, (U_i)_{i \in e}, (U_a)_{a \in \binom{e}{2}}, \ldots,
(U_a)_{a \in \binom{e}{d-1}}, U_e\right)\right)_{e \in \binom{\N}{d}},\]

\noindent for some middle-symmetric $f:[0,1]^{2^d} \to S$ and i.i.d. uniform random variables. Moreover, any array of this form is exchangeable.

\label{thm:representation_theorem_of_exchangeable_arrays}
\end{theorem}

\noindent A middle-symmetric function satisfies the following for any $\sigma \in \operatorname{Sym}([d])$
%
    \[f\left(x_\emptyset, (x_i)_{i \in [d]}, (x_a)_{a \in \binom{[d]}{2}}, \ldots, x_{[d]})\right) = f\left(x_\emptyset, (x_{\sigma(i)})_{i \in [d]}, (x_{\sigma(a)})_{a \in \binom{[d]}{2}}, \ldots,  x_{[d]})\right).\]
%
In words, the same permutation is applied coherently to the indices at every level.
More details can be found in~\cite{austin2013_exchangeable_random_arrays} or Chapter 7 of~\cite{kallenberg2005_probabilistic_symmetries}.

Throughout this chapter, we will encounter many domains of the form $[0,1]^{2^k}$ for some $k$. For ease of notation, we will assume the coordinates are always indexed by subsets of $[k]$ as in the definition of a middle-symmetric function above.

In the case of one-dimensional arrays, this representation theorem reduces to de Finetti's theorem, which says an exchangeable sequence is a mixture of i.i.d. sequences and thus takes the form $(X_n)\overset{d}{=}(f(U,U_n))$. Moreover, two-dimensional arrays take the form $f(U,U_i,U_j,U_{ij})$ for some function $f$ with $f(w,x,y,z)=f(w,y,x,z)$.

As \cite{aldous2010_more_uses_of_exchangeability} points out, this quantifies the fundamental relationship between
exchangeability and sampling that is already so well known: if we can get a hold of the representing function
$f$, it gives us a simple procedure for generating the corresponding combinatorial objects.
Next, we point out a simple way of generating new exchangeable arrays from old ones.

\begin{lemma}[State-Space Projections of Exchangeable Arrays]
    If $(X_e)$ is an exchangeable array with entries in $S$ and $g:S \to S'$ is measurable, then $(g(X_e))$ is an exchangeable array of the same dimension. Moreover, if $f$ is the representing function of $(X_e)$, then $g \circ f$ is the representing function of $(g(X_e))$.

    \label{lem:state_equivalence}
\end{lemma}

\noindent Here we are essentially defining an equivalence relation on the entries of the array where we let $s \sim s'$ if $g(s) = g(s')$.
We call this the \textit{state-space equivalence}.

Combined with Definition~\ref{def:exchangeable_array_process}, this gives three layers of Markov processes associated with an exchangeable array: (i) the initial labeled process corresponding to $(X_e)$, (ii) the labeled process corresponding to $(g(X_e))$, obtained by applying an equivalence to the state space, and (iii) the unlabeled chain, obtained by quotienting the coordinates. The ballot-sequence example below uses all three layers.


\begin{example}[P\'olya's Urn]
\label{ex:polyas_urn}
    In P\'olya's urn, an urn begins with $\alpha_i \in \N$ balls of color $i$ for $i \in [m]$. At each step,
    we sample a ball from the urn with replacement and add another ball of the same color as the sampled ball.
    It is well known that the sequence of balls drawn from a P\'olya urn is an exchangeable sequence with values in $[m]$,
    which is just a one-dimensional exchangeable array. In this case, the labeled chain $L_n$
    is the sequence of selected balls, and the unlabeled chain $M_n$ encodes how many balls of each color
    have been sampled. Classical theory implies this exchangeable array has representation function
    %
    \[f(x,y) = \sum_{i=1}^m i \cdot \ind \left\{y \in [t_{i-1}, t_{i})\right\},
    \qquad t_0:=0,\quad t_i := \sum_{k=1}^i p_k(x).\]
    %
    \noindent where $(p_1(U),\ldots,p_m(U)) \sim \mathrm{Dirichlet}(\alpha_1,\ldots,\alpha_m)$ for $\alpha_i$ initial balls of color $i$.
\end{example}

\begin{example}[Graphon Chain]
    Given any graphon $w$ there is a sequence of finite graphs $G_n$ converging to $w$ in the
    cut metric (equivalently the Doob-Martin topology by~\cite{grubel2015_persisting_randomness_graphs_search_trees})
    constructed via sampling $U_i \sim \operatorname{Unif}[0,1]$ and then connecting vertex $i$ to vertex
     $j$ with probability $w(U_i, U_j)$. Thus a graphon can be represented as an ``infinite adjacency matrix''
     which forms an exchangeable array with entries
    %
    \[X_e = \ind \{U_e \leq w(U_i, U_j)\} \quad \text{for } e = \{i,j\} \in \binom{\N}{2}.\]
    %
    \noindent Here the representation function $f(x_\emptyset, x_i, x_j, x_{ij})$ does not depend on $x_\emptyset$.
    \label{example:graphon_chain}
\end{example}

\begin{example}[d-Uniform Hypergraph Chain]
    We can extend the graphon example to $d$-uniform hypergraphs, which have hyperedges between subsets of $[n]$ of size $d$, initially explored in~\cite{elek2007_limits_hypergraphs_regularity_lemmas}. Namely, define a \textit{hypergraphon} as a middle-symmetric function $W:[0,1]^{2^d-2} \to [0,1]$ with coordinates indexed by proper, nonempty subsets of $[d]$. Then, analogous to graphons, this has representation function

    \[f\left(x_\emptyset, (x_i)_{i \in [d]}, \ldots, x_{[d]})\right) = \ind \Big \{x_{[d]} \leq W\left((x_i)_{i \in [d]}, \ldots,  (x_a)_{a \in \binom{[d]}{d-1}}\right) \Big \}.\]

    \noindent As we will see in Corollary~\ref{cor:classification_binary_array_processes} below, these are in a sense the only binary array processes.
\end{example}

\begin{example}[Ballot Sequences]
Given some total ordering of $I_0 := \bigcup_{i \in \N} \{a_i, b_i\}$ so $a_i < b_i$ for each $i$, there
are six possible ``interweavings'' of the letters
%
\[a_ib_ia_jb_j \quad a_i a_j b_i b_j \quad a_i a_j b_j b_i \quad a_j a_i b_i b_j \quad a_j a_i b_j b_i \quad a_j b_j a_i b_i.\]

\noindent This ordering can be encoded by a 2-dimensional array $X_{ij}$ with values in $\{1,\ldots,6\}$
describing which of these relationships the pairs $(a_i, b_i), (a_j, b_j)$ share.
If we equate the interweavings above which have the $a's$ and $b's$ in the same locations,
the unlabeled chain generates \textit{ballot sequences}: strings of $a's$ and $b's$ where every prefix
has at least as many $a's$ as $b's$. In~\cite{choi2023_doob_martin_growing_ballot_sequences}, the boundary is characterized through
a pair of diffuse marginals $(\zeta, \mu)$ encoding the relative density of $a's$ and $b's$ in the limit.
\end{example}

\noindent Exchangeable array processes also arise naturally in the theory of
exchangeable partitions through Kingman's paintbox theorem~\cite{austin2013_exchangeable_random_arrays}, random permutations through
permutons~\cite{grubel2023_ranks_copulas_permutons}, planar binary trees through R\'emy's chain~\cite{evans2014_doob_martin_remy_tree_growth}, and growing
random strings~\cite{choi2016_doob_martin_growing_random_words}. Many additional examples
of such combinatorial processes and their associated sampling representations can be found in
\cite{pitman2006_combinatorial_stochastic_processes}.

These and the preceding examples give a common strategy for embedding combinatorial processes as random arrays.
If a random structure consists of certain ``parts'' (e.g. vertices of a graph) and is completely determined by the
relationships within each subset of $d$ parts (e.g. which vertices are connected by an edge), then it can be
encoded as a $d$-dimensional array whose entries encode the relationship. Exchangeability then enforces the
constraint that the distribution of the structure is invariant under relabeling of the parts.

%%%%%%%%%%%%%%%%%%%%%%%%%%%%%%%%%%%%%%%%%%%%%
% Doob-Martin Boundary Theory for Exchangeable array Processes
%%%%%%%%%%%%%%%%%%%%%%%%%%%%%%%%%%%%%%%%%%%%%
\section{Doob-Martin Boundary of Exchangeable Array Processes}
\label{sec:exchangeable_array_boundary}

\noindent The backwards dynamics of the labeled chain are deterministic: given $L_n$ you simply restrict
to $e \in \binom{[n-1]}{d}$ to obtain $L_{n-1}$. Thus, its Doob-Martin boundary is almost trivial.
%
\begin{lemma}[Boundary of the Labeled Chain] \label{lem:boundary_labeled_chain}
    The boundary of $L_n$ consists of the deterministic arrays $x \in \mathrm{supp}(X)$.
\end{lemma}

\noindent Given the exchangeable array $(X_e)$ and a measurable function $g:S \to S'$, we can generate a new exchangeable
array $G(X)$ by applying $g$ coordinatewise: $G(X)_e := g(X_e)$. Although the boundary points of both labeled chains are deterministic arrays by Lemma~\ref{lem:boundary_labeled_chain},
state-space projection may collapse a nontrivial mixture of boundary points to a single projected boundary point. Indeed, any random array $X$ supported on $G^{-1}(y)$
trivially has $G(X) = y$ almost surely. Consequently, the bridge laws of the
original chain that project to y are precisely the probability measures supported on the fiber $G^{-1}(y)$.

%%%%%%%%%%%%%%%%%%%%%%%%
%%% Boundary Unlabeled Chain
%%%%%%%%%%%%%%%%%%%%%%%%
\subsection{Boundary of the Unlabeled Chain}

\noindent Now we seek to understand the boundary of the unlabeled chain $M_n$.
Recalling the correspondence \eqref{eqn:K_alpha_processes_char_h_processes} between infinite
 bridges of a Markov chain and its nonnegative harmonic functions via the Doob
 h-transform~\eqref{eqn:h-process} discussed in Appendix~\ref{appendix:doob_martin_background_appendix},
 we obtain the following lemma.

\begin{lemma}[Exchangeable Bridges of the Labeled Chain]
    Any infinite bridge of the labeled exchangeable array process is an array
    process, and it is exchangeable if and only if the corresponding harmonic
    function $h$ is constant on the equivalence classes of $\sim$.

    \label{lem:harmonic_functions_labeled}
\end{lemma}

\noindent This immediately implies the harmonic functions $\tilde{h}$ of the
unlabeled chain are in direct correspondence with the bridges of labeled chain which
come from exchangeable arrays.

\begin{corollary}[Bridges of the Unlabeled Chain]\label{cor:unlabeled_associated_to_exchangeable_array}
    If $B_n$ is an infinite bridge of the unlabeled chain $M_n$,
    then $B_n$ equals in distribution the unlabeled process associated to some
    exchangeable array $(X_e^h)$. In particular, if $\tilde{h}$ is the harmonic
    corresponding to the bridge $B_n$, then $h(x) := \tilde{h}([x])$
    is the harmonic corresponding to the array $(X_e^h)$.
\end{corollary}

The relation $\sim$ extends to infinite $d$-dimensional arrays $x$
in the obvious way: $x \sim y$ if there is $\sigma \in \operatorname{Sym}(\N)$ so $x = \sigma(y)$.
 We say a measurable collection $I$ of arrays is \textit{$\nu$-invariant} if
 permutations do not change it up to $\nu$-null sets
 %
    \[ I \overset{\nu-a.s.}{=} \tau I, \quad  \quad \forall \tau \in \operatorname{Sym}(\N).\]
%
Say an exchangeable array with distribution $\nu$ is \textit{ergodic} if $\nu(I) \in \{0,1\}$
for all $\nu$-invariant sets $I$, and say it is \textit{dissociated} if for any disjoint index
sets $I_1,\ldots,I_m \subset \N$ the arrays $\{X_e\}_{e \in \binom{I_i}{d}}$, $i\in[m]$, are independent.
For example, the exchangeable array corresponding to a graphon is dissociated: since edges are added independently,
the random subgraphs corresponding to disjoint collections of vertices are independent in the graphon model.
The following shows these properties characterize the minimal Doob--Martin boundary of
the unlabeled chain.

\begin{theorem}[Extremal Bridges and Ergodic Arrays]
    For any infinite bridge of the unlabeled exchangeable array process $M_n$, the following are equivalent:

    \begin{enumerate}
        \item The infinite bridge is extremal.
        \item Its corresponding exchangeable array is ergodic.
        \item Its corresponding exchangeable array is dissociated.
        \item Its corresponding exchangeable array admits a representation function $f$ that does not depend on $U_\emptyset$ almost surely.
    \end{enumerate}
    \label{thm:boundary_is_ergodic_arrays}
\end{theorem}


%
\begin{remark}
    By Example~\ref{example:graphon_chain}, given a graphon $w$ the representation function
    %
        \[f(q,x,y,u) := \ind \{u \leq w(x,y)\},\]
    %
    \noindent of the associated binary exchangeable array process does not depend on $q$.
    Hence, Theorem~\ref{thm:boundary_is_ergodic_arrays} implies
    graphons correspond to the extremal infinite bridges, and hence the minimal
    Doob-Martin boundary, of the graphon chain, as noted in~\cite{grubel2015_persisting_randomness_graphs_search_trees}.
\end{remark}

As a consequence, every unlabeled bridge admits a representation function
obtained by mixing the extremal representation
functions $f^{\tilde{K}(\cdot,\alpha)}$, which do not depend on the first
coordinate by the previous theorem. More specifically, if we
let $g: [0,1] \to \overline{E}$ be a function so that $g(U)$ has distribution
 $\mu_{\tilde{h}}$, we may choose
    %
    \[f^{\tilde{h}}(U, (U_i), \ldots, U_e) = f^{g(U)}((U_i), \ldots, U_e).\]
    %
\noindent Of course, the representation $f$ of a given exchangeable array
need not be unique. This is even apparent in the theory of graphons, which
are defined as isomorphism classes of functions $w:[0,1]^2 \to [0,1]$ modulo
Lebesgue isomorphisms $\lambda:[0,1] \to [0,1]$: we equate $w(\cdot, \cdot) \sim w(\lambda(\cdot), \lambda(\cdot))$.
The idea is that applying a Lebesgue isomorphism to a uniform does not change
its distribution, and thus will not influence the overall sampling distribution.
More generally, the representation functions of exchangeable arrays have a
similar uniqueness result ``up to Lebesgue-measure-preserving maps''.

The slight technicality is that we have to allow such Lebesgue-measure-preserving maps
to depend on the preceding coordinates. Namely, say a
map $T:[0,1]^{2^d} \to [0,1]$ \textit{preserves $\lambda$ in the last coordinate}
if for fixed $x_\emptyset, (x_i)_{i \in [d]}, \ldots, (x_e)_{e \in \binom{[d]}{d-1}}$ the map

\[y \mapsto T(x_\emptyset, (x_i)_{i \in [d]}, \ldots, (x_e)_{e \in \binom{[d]}{d-1}}, y),\]

\noindent preserves the Lebesgue measure. For a proof of the following, see
Theorem 7.28 in~\cite{kallenberg2005_probabilistic_symmetries}, which proves
 a more general statement that does not restrict the array to a finite dimension $d$.
 For something easier to parse, the case $d=2$ can also be found in
 Theorem 2.1' of \cite{kallenberg1989_representation_exchangeable_arrays}.
%
\begin{theorem}[Uniqueness of Representation Functions]
    The following are equivalent:
    \begin{enumerate}
        \item $f,g:[0,1]^{2^d} \to S$ generate the same exchangeable array (in distribution).

        \item There exists measurable functions $T_i,T_i' : [0,1]^{2^i} \to [0,1]$  for  $i = 0,\ldots, d$ which preserve $\lambda$ in the last coordinate so that a.s.
        %
        \begin{align*}
          f\Big( & T_0(x_\emptyset), (T_1(x_\emptyset, x_i))_{i \in [d]}, \dots, T_d(x_\emptyset, (x_i)_{i \in [d]}, \dots, x_{[d]}) \Big)  \\
          & = g\Big(  T'_0(x_\emptyset), (T'_1(x_\emptyset, x_i))_{i \in [d]}, \dots,  T'_d(x_\emptyset, (x_i)_{i \in [d]}, \dots, x_{[d]}) \Big).
        \end{align*}
    \end{enumerate}
    \label{thm:uniqueness_of_rep_functions}
\end{theorem}
%
In the $d=2$ case, this exactly recovers the uniqueness statement for graphons discussed above.
An immediate corollary of Theorem~\ref{thm:uniqueness_of_rep_functions} is that, in a sense, the $d$-uniform hypergraph chain mentioned in the previous examples is the \textit{only} binary exchangeable array process.

\begin{corollary}[Binary Representation Functions]
    Any binary exchangeable array admits a (not necessarily unique) representation function of the form
    %
    \[f = \ind \Big\{x_{[d]} \leq w\left(x_\emptyset, (x_i)_{i \in [d]}, \ldots, (x_e)_{e \in \binom{[d]}{d-1}} \right)\Big\},\]
    %
    \noindent for some $w:[0,1]^{2^d - 1} \to [0,1]$.
    \label{cor:classification_binary_array_processes}
\end{corollary}

Again, we are primarily interested in the extremal representation functions, which
depend trivially on $x_\emptyset$ by Theorem~\ref{thm:boundary_is_ergodic_arrays}.
Hence our hypergraphons can be viewed as functions
    %
    \[w:[0,1]^{2^d-2} \to [0,1],\]
    %
\noindent indexed by proper nonempty subsets of $[d]$. This recovers the
original notion introduced previously in~\cite{elek2012_measure_theoretic_dense_hypergraphs} and whose relationship to
exchangeability was further explored in~\cite{austin2008_exchangeable_random_variables}.


%%%%%%%%%%%%%%%%%%%%%%%%%%%%%%%%%%%%%%%%%%%%%%%%%%%%%%
% Embedding Density of Representation Functions
%%%%%%%%%%%%%%%%%%%%%%%%%%%%%%%%%%%%%%%%%%%%%%%%%%%%%%
\section{Embedding Density of the Representation Functions}
\label{sec:embedding_density_representation_functions}

\noindent The idea of embedding a smaller structure into a larger one plays a pivotal
role in the theory of combinatorial processes. Two of our principal examples,
graphons and permutons, already come equipped with natural notions of
embedding: subgraphs and subpermutations respectively. Somehow these embeddings encode the ``local'' structure of the larger combinatorial
object. We can then compare two such objects by comparing these local structures at various
different scales.


The aim of this section is to recast these notions in the language
of exchangeable arrays, with particular emphasis on binary arrays, and to
extend the framework developed in~\cite{borgs2007_convergent_dense_graph_sequences}. For a finite binary $d$-dimensional
array $F$ of size $k$, let
%
\[\ind(F) := \left\{S \in \binom{[k]}{d} : F_S = 1\right\}.\]
%
\noindent denote the collection of entries of our array which are one.
Moreover, given a vector $x := (x_e)_{e \subseteq [n]} \in [0,1]^{2^n}$
and a subset $A \in \binom{[n]}{d}$ denote
%
\[x_{(A)} := \left(x_\emptyset, (x_i)_{i \in A}, (x_{ij})_{ij \in \binom{A}{2}}, \;\ldots \;, x_A\right) \in [0,1]^{2^d}.\]
%
\noindent Namely,
if $x$ is indexed by all subsets of $[n]$ then $x_{(A)}$ restricts to
only subsets of $A$. Lastly, for a finite $d$-dimensional array of size
$k$ let $|x| := k$ denote its size.

%%%%%%%%%%%%
% Embeddings of Labeled States
%%%%%%%%%%%%

\subsection{Embeddings of Labeled States}

Suppose we have arrays $(x_e)_{e \in \binom{[n]}{d}}$ and $(y_e)_{e \in \binom{[n+m]}{d}}$.
Then we say \textit{x is embedded in y} if $y_e = x_e$ for
all $e \in \binom{[n]}{d}$ and denote this by $x \leq y$. Namely, we obtain $x_e$ by
looking at the ``top left corner'' of the array $y_e$. From the viewpoint of our
labeled Markov chain $L_n$, this just means that the array $(y_e)$ is reachable starting in
state $(x_e)$. This extends the standard notion of an induced subgraph beyond binary arrays.

We say a binary array $x$ is \textit{weakly embedded} in binary array $y$, denoting this
by $x \lesssim y$, if $y_e = 1$ whenever $x_e = 1$. This is the standard notion of a graph embedding.

%%%%%%%%%%%%%%%%%%%%%
% Embeddings of Unlabeled States
%%%%%%%%%%%%%%%%%%%%%
\subsection{Embeddings of Unlabeled States}

\noindent Again suppose we have arrays $(x_e)_{e \in \binom{[n]}{d}}$ and $(y_e)_{e \in \binom{[n+m]}{d}}$.
Say $[x]_\sim$ is embedded in $[y]_\sim$ if there exists $\sigma \in \operatorname{Sym}([n+m])$ so $(x_e) \leq (y_{\sigma(e)})$.
Equivalently, there exists an injective function $\phi:[n] \to [n+m]$ so $(x_{e}) = (y_{\phi(e)})$.

In this case, there can be more than one embedding if there are several
different choices of $\sigma$ that produce an embedding.
This suggests a notion of \textit{embedding density} analogous to the
theory of graphons and permutons.
%
\begin{align}
    emb(x,y) & := \frac{| \{\sigma \in \operatorname{Sym}([n+m]) \; : \; (x_e) \leq (y_{\sigma(e)}) \}|}{(n+m)!} \notag\\[7pt]
            & = \frac{| \{\phi:[n] \xrightarrow{injective} [n+m] \; : \; (x_{e}) \leq (y_{\phi(e)}) \}|}{(n+m) (n+m-1) \cdot \ldots \cdot (m+1)}
    \label{def:embedding_density}
\end{align}

\noindent for $x \in E_n, y \in E_{n+m}$. Note that the embedding
density is constant on each equivalence class of $x$ so this is
well-defined as a function on these equivalence classes.
For binary arrays, we analogously define the \textit{weak embedding density}
$emb^W$ by replacing $\leq$ with $\lesssim$. The exact and weak embedding
densities satisfy the following relationships.

\begin{lemma}[Exact and Weak Embedding Densities]
    For binary arrays $x,y$ we have

    \begin{enumerate}[label=(\roman*)]
        \item $emb(x,y) \leq emb^W(x,y)$.
        \item $emb^W(x,y) = \sum_{x' \geq x} emb(x',y)$.
        \item $emb(x,y) = \sum_{x' \geq x} (-1)^{|\ind(x')| - |\ind(x)|} \cdot emb^W(x',y)$.
    \end{enumerate}

    \label{lem:relating_different_notions_of_embedding}
\end{lemma}

\noindent Given any notion $t(x,y)$ of the embedding density of $x$ into $y$, we can topologize our collection
of combinatorial objects: we simply say a sequence $y_n$ converges if
%
    \[\lim_{n \to \infty} t(x, y_n) \text{ exists}\]
%
for every fixed $x$. Thus Lemma~\ref{lem:relating_different_notions_of_embedding}
shows that exact and weak embedding densities generate the same topology for
binary arrays.
Such convergence and its associated topology is often referred to as \textit{left-convergence}
in the graphon literature. In fact, left-convergence is closely related to the Doob-Martin topology.

\begin{lemma}[Doob-Martin Convergence and Embedding Densities]
    Let $y_N\in E_N$. The following are equivalent:
        \begin{enumerate}[label=(\roman*)]
            \item The sequence $y_N$ converges in the Doob-Martin topology.
            \item The embedding density $emb(x, y_N)$ converges for all $x$.
        \end{enumerate}

    \noindent For binary arrays, convergence of the weak embedding density is also equivalent.

    \label{lem:equivalence_of_dm_and_density_convergence}
\end{lemma}

\noindent In the case of our graphon and permuton chains, the notion of
embedding can be extended to the limiting object, graphons and permutons,
 via sampling. Namely, for a graphon $w$ its embedding density $t(\cdot, w)$
  is essentially defined so that for each $n$ and finite graph $F$
%
\[t(F, w) = \E t(F, G_n(w)),\]

\noindent where $G_n(w)$ is the random graph generated by sampling
 $n$ uniforms $U_1,\ldots, U_n$ and connecting vertices $i,j$ with
 probability $w(U_i, U_j)$.

Analogously, we would like to extend this notion of an embedding density
to our limiting objects: the representation functions $f$ of our
exchangeable arrays. This suggests
%
\begin{equation}
    emb(x, f) := \E \; emb(x, X_n(f)), \qquad  emb^W(x, f) := \E \; emb^W(x, X_n(f)),
    \label{eqn:def_embedding_density_representation_function_f}
\end{equation}

\noindent where $X_n(f)$ is the array indexed by $\binom{[n]}{d}$ generated by
%
\[(X_n(f))_e := f(U_\emptyset, (U_i)_{i \in e},\ldots, U_e).\]
%
\noindent We will now show these expectations are independent of $n$ and
have an integral representation analogous to the embedding density of graphons.

\begin{lemma}[Embedding Densities of Representation Functions]
Suppose $d \leq k \leq n$. If $F$ is a size-$k$, $d$-dimensional array and
$f$ represents an exchangeable array $X$, then, independently of $n$,

\begin{equation}
    emb(F,f) = \int_{[0,1]^{2^k}}  \prod_{S \in \binom{[k]}{d}} \ind \{f(x_{(S)}) = F_S\} \; d\lambda_{2^k}.
\end{equation}
If $F$ and $f$ are binary-valued, then also
\begin{equation}
    emb^W(F,f) = \int_{[0,1]^{2^k}}  \prod_{S \in \ind(F)} f(x_{(S)}) \; d\lambda_{2^k}.
    \label{eqn:integral_formula_for_embedding_density_of_representing_function}
\end{equation}
%
\noindent Moreover, if $F$ and $f$ are binary-valued and $f$ does not depend on $x_\emptyset$, we have
%
\begin{align*}
    \P\left( |emb^W(F, X_n(f)) - emb^W(F, f)| > \epsilon \right) &  \leq 2 \exp \left(-\frac{\epsilon^2}{2k^2} n \right),\\
    %
    \P\left( |emb(F, X_n(f)) - emb(F, f)| > \epsilon \right) & \leq 2 \exp \left(-\frac{\epsilon^2}{2k^2} n \right).
\end{align*}
%
\noindent Thus,
\[
\lim_{n \to \infty} emb(F, X_n(f)) \overset{a.s.}{=} emb(F,f),
\]
\noindent and
\[
    \lim_{n \to \infty} emb^W(F, X_n(f)) \overset{a.s.}{=} emb^W(F,f).
\]
\noindent In particular, the results of Lemma~\ref{lem:relating_different_notions_of_embedding} (i)--(iii) extend to $emb^W(F,f)$ and $emb(F,f)$.


\label{lem:integral_representations_of_embedding_densities}
\end{lemma}

\noindent This is analogous to the graphon identity
%
\[\E\,emb^W(F,G_n(w))
= \int_{[0,1]^k} \prod_{\{i,j\} \in \ind(F)} w(x_i,x_j)\,d\lambda_k(x).\]
%
\noindent To see why, note when $d=2$
%
\[ \int_{[0,1]^{2^k}}  \prod_{S \in \ind(F)} f(x_{(S)}) \; d\lambda_{2^k} = \int_{[0,1]^{2^k}}  \prod_{\{i,j\} \in \ind(F)} f(x_\emptyset, x_i, x_j, x_{ij}) \; d\lambda_{2^k}. \]
%
\noindent For graphons the terms inside the product become $\ind\{x_{ij} \leq w(x_i, x_j)\}$ so
we can integrate out the coordinate $x_\emptyset$ and the coordinates $x_{ij}$ (which each only appear
 in one term of the product) to get the familiar formula.

In general, we can see that the integral in
\eqref{eqn:integral_formula_for_embedding_density_of_representing_function} only relies
on the $\sum_{i=0}^d \binom{k}{i}$ coordinates corresponding to subsets $e \subseteq [k]$ of size at most $d$,
so many of the coordinates in the integral are superfluous. We leave them in for simplicity of notation.

Lastly, we will show one strategy for embedding a \textit{finite} binary
exchangeable array $z$ into the space of representation functions. A natural candidate
 is just a continuous version of the hypergraph adjacency matrix,
 analogous to how we associate a graphon to a finite graph
    %
    \begin{equation}
        f_z(u_\emptyset, (u_i)_{i \in [d]}, \ldots, u_{[d]})
        := \ind \left\{\{\lceil |z|u_i\rceil:i\in[d]\} \in \ind(z) \right\}.
        \label{eqn:ex_array_sampling_function}
    \end{equation}

    %
\noindent We call $f_z$ the \textit{sampling function} associated to a finite array
$z$. Sampling from $f_z$ chooses vertices of $z$ independently and hence with
replacement, whereas the embedding density samples without replacement. The
resulting discrepancy is controlled by the probability of a collision, and
vanishes asymptotically.

\begin{lemma}[Sampling Functions Approximate Weak Embedding Densities]\label{lem:sampling_functions_approximate_embedding_densities}
    Let $F$ and $z$ be finite binary $d$-dimensional arrays with
    $|F|\leq |z|$. Then
    \[
        \left|emb^W(F,f_z)-emb^W(F,z)\right|
        \leq \frac{1}{|z|}\binom{|F|}{2}.
    \]
\end{lemma}



%%%%%%%%%%%%%%%%%%%%%%%%%%%%%%%%%%%%%%%%%%
% Conclusions and Future Directions
%%%%%%%%%%%%%%%%%%%%%%%%%%%%%%%%%%%%%%%%%%
\section{Conclusions and Future Directions}

This chapter provides a framework for viewing a broad class of combinatorial
growth processes through exchangeable arrays. Using the classical Aldous--Hoover--Kallenberg
representation theorem for such arrays, we can identify such arrays with sampling functions
on some hypercube, generalizing the familiar graphon picture. Applying the Doob-Martin boundary
theory, we show that the points of the minimal boundary can equivalently be viewed as
ergodic arrays or through representation functions that do not depend on the first
coordinate $x_\emptyset$.

We define exact and weak notions of embedding one finite array into another
and show that they generate the same topology for binary arrays. Using this,
we show that this probabilistic viewpoint of the limit objects is equivalent to the
more combinatorial notion of left-convergence built from these embedding densities, readily generalizing the graphon literature.
In this way, we extend the relationship between the classical Doob-Martin boundary and these embedding
densities previously explored in~\cite{grubel2015_persisting_randomness_graphs_search_trees,grubel2023_ranks_copulas_permutons}.
This identifies a common compactified limit space, a feature that is far from obvious in the original
construction of~\cite{lovasz2004_limits_dense_graph_sequences}.

Our work so far omits one key success of graphon theory: the cut metric. For graphons, this metric metrizes
the topology of left-convergence. It has a direct combinatorial interpretation, controlling discrepancies
in edge densities across all pairs of vertex subsets and, through the counting lemma, the densities of every
fixed finite graph. In higher dimensions and for nonbinary arrays, however, the analogous metric theory appears more delicate.
In Appendix~\ref{appendix:cut_metric_partial_progress}, we record some preliminary directions. In particular, we
discuss a natural extension to hypergraphs and prove a corresponding hypergraph counting lemma. We then give a
random sequence that left-converges in probability while its canonical representatives fail to be Cauchy in probability
in the resulting strong cut norm. Furthermore, we discuss some ways one might be able
to extend any such metric from these hypergraphs (e.g. binary exchangeable arrays) to exchangeable arrays with more general
state spaces via a one-hot encoding.




%%%%%%%%%%%%%%%%%%%%%%%%%%%%%%%%%%%%%%%%%%
% Metrizing Convergence of Exchangeable Array Processes %%%%%%%%%%%%%%
%%%%%%%%%%%%%%%%%%%%%%%%%%%%%%%%%%%%%%%%%%
\newpage
\begin{subappendices}

\section{Candidate Cut Metrics for Exchangeable Array Processes}
\label{appendix:cut_metric_partial_progress}

As Lemma~\ref{lem:equivalence_of_dm_and_density_convergence} and~\cite{grubel2015_persisting_randomness_graphs_search_trees}
show, graphons and the Doob-Martin theory are topologically equivalent.
While the latter already comes equipped with a metric, this
metric is not naturally interpretable in terms of the underlying combinatorial process.
The real success of the theory of graphons was the introduction of an alternative metric, the cut metric, which
is reasonably interpretable, explicitly encodes graph-theoretic properties, and allows for exact
computation of many quantities of interest.

This appendix records preliminary progress towards a cut-type metric for
exchangeable array processes. First, we discuss the binary case, where
$d$-dimensional arrays are closely related to $d$-uniform hypergraph chains.
Then we describe a possible extension to nonbinary arrays by projecting the
state space onto a collection of binary arrays.

We saw in the above example that the graphon $w$ of a graphon
chain is closely related to the representation function $f$ of
its associated exchangeable array. Namely,
 %
\[f(q,x,y,u) = \ind \{u \leq w(x,y)\}.\]
 %
\noindent Hence we hope to reinterpret the cut metric in terms of a metric
on these representation functions $f$, and from there find a way to extend
this metric to arbitrary exchangeable array processes. For fixed graphon
representatives, recall the cut norm distance
%
\begin{align*}
    \lVert w-w'\rVert_\Box & := \sup_{0 \leq g,h \leq 1} \left | \int_{[0,1]^2}
 g(x)h(y)(w(x,y) - w'(x,y)) \; dx \; dy\right|.
\end{align*}
%
\noindent This suggests a natural generalization to exchangeable array processes.
Given a hypergraphon $W:[0,1]^{2^d-2}\to[0,1]$ associated with a
$d$-dimensional binary exchangeable array, define

\[
\lVert W\rVert_{\Box^{d-1}}
:= \sup_{g_1,\ldots,g_d}
\left|\int_{[0,1]^{2^d-2}} W(x)
\prod_{i=1}^d g_i\left((x_A)_{\emptyset\neq A\subseteq[d]\setminus\{i\}}\right)
\,d\lambda_{2^d-2}(x)\right|,
\]
where the supremum is over measurable functions
$g_i:[0,1]^{2^{d-1}-1}\to[0,1]$.

\noindent While this form may seem
strange at first, it arises naturally from attempting to extend the classical counting
lemma to hypergraphons.

\begin{lemma}[Counting Lemma] \label{lem:counting_lemma}
    Let $F$ be a finite binary $d$-array of size $k\geq d$, and let $W,W'$ be
    $d$-uniform hypergraphons. Then
    \[
    |emb^W(F,W)-emb^W(F,W')|
    \leq |\ind(F)|\,\lVert W-W'\rVert_{\Box^{d-1}}.
    \]
\end{lemma}

\begin{proof}
    Write $x_{(<S)}:=(x_A)_{\emptyset\neq A\subsetneq S}$. By the integral
    representation of the hypergraphon embedding density,
    \begin{align*}
        \text{LHS} &= \left | \int_{[0,1]^{2^k}}
        \prod_{S \in \ind(F)} W(x_{(<S)})
        - \prod_{S \in \ind(F)} W'(x_{(<S)}) \, d\lambda_{2^k} \right |
        %
        \intertext{Enumerate $\ind(F)=\{S_1,S_2,\ldots,S_\ell\}$. Interpolating between the products one factor at a time and applying the triangle inequality gives}
        &\leq \sum_{m=1}^\ell \left| \int_{[0,1]^{2^k}}
        \left( \prod_{i=1}^{m-1} W(x_{(<S_i)})\right)
        (W(x_{(<S_m)})-W'(x_{(<S_m)}))
        \left( \prod_{i=m+1}^{\ell} W'(x_{(<S_i)})\right)
        \,d\lambda_{2^k} \right|.
    \end{align*}
    For each summand, condition on all coordinates not indexed by nonempty
    proper subsets of $S_m$. Grouping the remaining factors according to their
    intersections with $S_m$ puts the inner integral in the form defining
    $\lVert W-W'\rVert_{\Box^{d-1}}$. Each summand is therefore at most this
    cut norm, and summing over the $\ell=|\ind(F)|$ factors proves the result.
\end{proof}

As in the case of graphons, a hypergraphon representation is not unique, so
we minimize over possible representations
%
\[d_\Box(W,W') := \inf_{V \sim W,\,V' \sim W'} \lVert V-V'\rVert_{\Box^{d-1}}.\]
%
\noindent Here $V\sim W$ means that $V$ and $W$ generate the same exchangeable-array law.
%
\noindent The counting lemma shows that convergence in the quotient distance $d_\Box$ implies convergence
in the topology generated by embedding densities, and hence also in the Doob-Martin topology by
Lemma~\ref{lem:equivalence_of_dm_and_density_convergence}. The converse fails
for the strong cut norm on a fixed choice of representatives: left-convergence
need not make the canonical step representatives Cauchy in this norm. The
following example, adapted from~\cite{zhao2015_hypergraph_limits_regularity_approach}, illustrates this distinction.

\begin{example}
    Let $(G_n)_{n\geq1}$ be independent Erd\H{o}s--R\'enyi graphs with
    $G_n\sim G(n,1/2)$. Define a sequence of three-dimensional hypergraphs
    $H_n$ by taking the triangles of $G_n$ as hyperedges. Then, for any fixed
    three-dimensional hypergraph $F$,
    %
        \[emb^W(F,H_n) \xrightarrow{p} 2^{-|\partial F|},\]
    %
    \noindent where $\partial F$ is the set of unordered pairs of vertices
    contained in some hyperedge of $F$. Hence $H_n$ left-converges in probability to the
    hypergraphon $W(x)=\ind\{x_{12},x_{13},x_{23}\leq1/2\}$. Let
    $W_n$ be the step hypergraphon associated with $H_n$. Then
    %
        \[W_n(u_1,u_2,u_3,u_{12},u_{13},u_{23})
        = A_n(u_1,u_2)A_n(u_1,u_3)A_n(u_2,u_3),\]
    %
    where $A_n:[0,1]^2\to\{0,1\}$ is the graphon associated
    to $G_n$, namely
    %
        \[A_n(x,y) = \ind\{(\lceil nx \rceil, \lceil ny \rceil) \in E(G_n)\}.\]
    %
    \noindent Hence, for $m>n$,
    %
    \begin{align*}
        \lVert W_n-W_m\rVert_{\Box^2}
        &\geq \left|\int_{[0,1]^3}(W_n-W_m)
        A_n(x,y)A_n(x,z)A_n(y,z)\,dx\,dy\,dz\right|.
    \end{align*}
    %
    \noindent Note
    %
        \[\int_{[0,1]^3}W_n A_n(x,y)A_n(x,z)A_n(y,z)\,dx\,dy\,dz
        = emb^W(K_3,G_n) \xrightarrow{p} \frac18.\]
    %
    \noindent Similarly, independence of $G_n$ and $G_m$ gives, as $n,m\to\infty$,
    %
        \[\int_{[0,1]^3}W_m A_n(x,y)A_n(x,z)A_n(y,z)\,dx\,dy\,dz
        = \frac{1}{64}+o_p(1),\]
    %
    \noindent since, outside a set of measure $O(n^{-1}+m^{-1})$, the integrand
    checks for three distinct edges in each of the two independent graphs. Hence
    $\lVert W_n-W_m\rVert_{\Box^2}\geq 7/64+o_p(1)$, so these canonical
    representatives are not Cauchy in probability in $\lVert\cdot\rVert_{\Box^2}$.
\end{example}

\noindent This calculation concerns the displayed representatives $W_n$ and
does not by itself rule out convergence after minimizing over equivalent
representations in the quotient distance $d_\Box$.

\noindent The paper~\cite{zhao2015_hypergraph_limits_regularity_approach} attempts to circumvent this issue by
replacing convergence in the strong cut norm with a weaker notion of \textit{partitionable convergence},
based on compatible hierarchies of weak regularity partitions of lower-order faces. This notion admits subsequential
 limits and implies convergence of homomorphism densities. Since both the Doob-Martin and left-convergence constructions
 yield compact limit spaces, this is one step in the right direction. However, \cite{zhao2015_hypergraph_limits_regularity_approach}
 still leaves the converse conjectural.


% Cut metric for Nonbinary Arrays %%%%%%%%%%%%%%%%%%%%%%%%%%%%%%%%%%%%%%%%%%%
\newpage
\subsection{Cut metric for nonbinary arrays}

In this section, we provide partial progress towards extending the cut norm beyond hypergraphons
to nonbinary arrays. The idea is to characterize a nonbinary array in terms of a series of projections onto
the space of binary arrays, and apply the existing cut metric to each projection separately.

First, we need to extend the notion of a weak embedding to the nonbinary case. For a subset of states
 $C \subseteq S$ define a \textit{$C$-weak embedding} of $x$ into $y$ to be an injective map
 $\phi: [|x|] \to [|y|]$ which preserves the states in $C$. Namely, if $s \in C$ then $x_e = s$
 implies $y_{\phi(e)} = s$. In the case $C = \{1\}$ and $S = \{0,1\}$ this reduces to the standard
  notion of weak embeddings of a binary array. Denote the set of such embeddings by
  $\operatorname{Emb}_C(x,y)$ and their density by $emb^C(x,y)$. Then we have the following extension
  of Lemmas~\ref{lem:relating_different_notions_of_embedding} and
  \ref{lem:equivalence_of_dm_and_density_convergence}.

\begin{lemma}[Weak Embeddings for Nonbinary Arrays]
    For finite $S$ and any $C \subseteq S$:

    \begin{enumerate}[label=(\roman*)]
        \item $emb^C(x,y) = \sum_{x' \geq_C x} emb(x', y)$.

        \item $emb(x,y) \leq emb^C(x,y)$.
    \end{enumerate}

    \noindent Moreover, if $C = S \setminus \{s\}$ then:

    \begin{enumerate}
        \item[(iii)] $emb(x,y) = \sum_{x' \geq_C x} (-1)^{N_s(x)-N_s(x')} \cdot emb^C(x', y)$.
    \end{enumerate}


    \noindent where $x' \geq_C x$ means $x'$ can be obtained from $x$ by changing some of the values $x_e$ with $x_e \notin C$ and $N_s(x) := | \{e \; : \; x_e = s\} |$ is the number of entries of $x$ which are $s$.

    In particular, this shows that convergence of $emb(\cdot, y_N)$ implies convergence of $emb^C(\cdot, y_N)$ for any $C$, and the converse is true if $C = S \setminus \{s\}$.

    \label{lem:weak_embeddings_for_nonbinary_arrays}
\end{lemma}

\begin{proof}
Write $\operatorname{Emb}(x,y)$ and $\operatorname{Emb}_C(x,y)$ for the
sets of embeddings and $C$-weak embeddings, respectively.

For (i), let $\phi$ be a $C$-weak embedding of $x$ into $y$. There is
exactly one array $x' \geq_C x$ for which $\phi$ is an embedding of $x'$
into $y$: set $x'_e=x_e$ when $x_e\in C$, and set
$x'_e=y_{\phi(e)}$ when $x_e\notin C$. Thus the set of $C$-weak
embeddings is the disjoint union of the sets of embeddings of
$x' \geq_C x$ into $y$. Dividing by the number of injective maps
$[|x|]\to[|y|]$ proves (i).

For (ii), every embedding of $x$ into $y$ is automatically a $C$-weak
embedding. Consequently, the embedding set is contained in the set of
$C$-weak embeddings. Dividing their cardinalities by the common number of
injective maps gives
\[
    emb(x,y) \leq emb^C(x,y).
\]

For (iii), suppose $C=S\setminus\{s\}$ and let
$E_s(x)=\{e:x_e=s\}$. For $e\in E_s(x)$ and $t\in C$, let $A_{e,t}$ be
the set of $C$-weak embeddings $\phi$ for which $y_{\phi(e)}=t$. An exact
embedding of $x$ is a $C$-weak embedding that belongs to none of these
sets. Hence inclusion--exclusion gives
\[
    |\operatorname{Emb}(x,y)|
    = \sum_{x'\geq_C x}
      (-1)^{N_s(x)-N_s(x')}
      |\operatorname{Emb}_C(x',y)|.
\]
Indeed, choosing an intersection of the events $A_{e,t}$ amounts precisely
to choosing the positions at which an entry $s$ of $x$ is replaced, along
with its replacement state $t\in C$; the resulting array is the
corresponding $x'\geq_C x$. Dividing by the common number of injective maps
proves (iii).
\end{proof}


\noindent There is a simple way of encoding a nonbinary array with finite
state space $S$ as a collection of binary arrays of the same dimension:
simply apply a one-hot encoding to the state space. More specifically,
for each $A \subseteq S$ define the function $g_A: S \to \{0,1\}$ by
$g_A(s) = \ind \{s \in A\}$. Thus each $d$-dimensional array with values
in $S$ is encoded as $2^{|S|}$ binary arrays of the same size. First,
we show there is a relationship between embeddings in the original
state space and embeddings of the corresponding binary arrays.

% Lemma %
\begin{lemma}[Binary Projections of Nonbinary Arrays]
    For $d$-dimensional arrays $x,y$ with values in some finite set $S$,
    %
        \[\operatorname{Emb}(x,y)
        = \bigcap_{s \in S}\operatorname{Emb}(g_{\{s\}}(x),g_{\{s\}}(y)).\]
    %
    \noindent Moreover, if $y_N$ left-converges to $f$, then $g_A(y_N)$ left
    converges to $g_A \circ f$ for all $A \subseteq S$.
    \label{lem:reducing_nonbinary_conv_to_binary}
\end{lemma}

\begin{proof}
% proof formula
    Suppose $\phi:[|x|] \to [|y|]$ is injective. Then $x_e=y_{\phi(e)}$
    if and only if $g_{\{s\}}(x_e)=g_{\{s\}}(y_{\phi(e)})$ for every
    $s\in S$, proving the first statement.

    For the second, fix a binary array $b$ and let
    \[
        \mathcal{Y}_A^b:=\left\{x\in S^{\binom{[|b|]}d}:
        b_e=1\Rightarrow x_e\in A\right\}.
    \]
    Every weak embedding of $b$ into $g_A(y_N)$ induces a unique exact array
    $x\in\mathcal{Y}_A^b$, so
    \[
        emb^W(b,g_A(y_N))=\sum_{x\in\mathcal{Y}_A^b}emb(x,y_N).
    \]
    Since $S$ is finite and $y_N$ left-converges to $f$, this converges to
    \[
        \sum_{x\in\mathcal{Y}_A^b}emb(x,f)=emb^W(b,g_A\circ f).
    \]
\end{proof}

\noindent We conjecture that the converse is true as well.

\begin{conjecture}
    If $g_A(y_N)$ left-converges to $g_A\circ f$ for all $A \subseteq S$, then $y_N$ left-converges to $f$.
\end{conjecture}

Given an exchangeable array process $(X_e)$ represented by $f$,
Lemma~\ref{lem:state_equivalence} implies the array process $(g_A(X_e))$ is
a binary exchangeable array process represented by $g_A \circ f$. Thus if $f,f'$ share
the same state space $S$, if the given conjecture is true then it is natural to define
%
\begin{equation}
    d(f,f') := \frac{1}{2^{|S|}} \sum_{A \subseteq S} d(g_A \circ f, \; g_A \circ f').
\end{equation}
%




%%%%%%%%%%%%%%%%%%%%%%%%%%%%%%%%%%%%%%%%%%%%%%%%%%%%%%%%%%%%%%%%%%%%%%
% Doob-Martin Background Appendix
%%%%%%%%%%%%%%%%%%%%%%%%%%%%%%%%%%%%%%%%%%%%%%%%%%%%%%%%%%%%%%%%%%%%%%
\newpage
\section{Background on Doob-Martin Boundary Theory}
\label{appendix:doob_martin_background_appendix}
Here we aim to provide some of the essentials of Doob-Martin boundary theory.
We restrict our attention to the case of discrete Markov chains relevant to this
chapter. The primary reference on this topic is \citep{doob1959_discrete_potential_theory_boundaries} but
\citep{sawyer_martin_boundaries_random_walks} and \citep{evans2014_doob_martin_remy_tree_growth} also provide good overviews for the discrete
case in more modern terminology.\\

 Given a transient Markov chain $\Xn$ with a countable state space $E$ along with a base state $x_0 \in E$ from which the chain is irreducible, some of the goals of Martin boundary theory are:

\begin{enumerate}[label=(\roman*)]
    \item Complete $E$ into a compact metric space by adding ``points at $\infty$'' to which the Markov chain is converging.
    %
    \item Characterize all the nonnegative/bounded harmonic functions $h :E \to [0,\infty)$.
    %
    \item Develop a notion of conditioning the chain ``at infinity'' and understanding the behavior of these conditioned chains.
\end{enumerate}

%
\indent Suppose $\Xn$ has transition kernel $p$. Define the \textit{Green's kernel} or
\textit{potential kernel} $g: E \times E \to \mathbb{R}_{\geq0}$ by
%
\begin{equation}
    g(x,y) := \sum_{n \geq 0} p^n(x,y) = \E_x (\text{number of times $X_n$ hits y}).
    \label{eqn:green's_kernel}
\end{equation}
%
\noindent Transience implies $g(x,y) < \infty$. Assume further there exists a base state $x_0$ from which every state is
accessible, namely $g(x_0,y) > 0$ for all $y \in E$. For this chapter, $x_0$ will
typically be the empty array.
Moreover, we will typically be interested in \textit{graded} chains. These are
Markov chains in which the state space $E  =\bigsqcup E_n$ can be partitioned
so transitions occur only from $E_n$ to $E_{n+1}$. Here
%
    \[g(x,y) = \P_x(X_n \text{ hits }y).\]
%
\noindent Define the \textit{Doob-Martin kernel} $K:E \times E \to \R_{\geq 0}$
with respect to base state $x_0$ via
%
\begin{equation}
    K(x,y) := \frac{g(x,y)}{g(x_0,y)} = \frac{\P_x(X_n \text{ hits y})}{\P_{x_0}(X_n \text{ hits y})}.
\end{equation}

%
\noindent Observe for each fixed $y$ the function $K(\cdot, y)$ is ``almost harmonic'' in the sense

\begin{equation*}
    \sum_{z \in E} p(x,z)K(z, y) = K(x,y) \quad x \neq y.
\end{equation*}


\noindent This suggests if we let $y$ ``go off to infinity'' we ought to get a truly
harmonic function in the limit. The Doob-Martin boundary will make this intuition precise.

\indent Note for fixed $x$, the kernel $K(x, y)$ is a bounded function of $y$. To see why,
 note accessibility implies $p^m(x_0, x) > 0$ for some $m$, and the Markov
property gives $g(x_0,y)\geq p^m(x_0,x)g(x,y)$. Thus
%
if $g(x,y)=0$, then $K(x,y)=0$; otherwise,
\[0 \leq K(x,y) = \frac{g(x,y)}{g(x_0, y)}
\leq \frac{1}{p^m(x_0,x)} =: C_x.\]

%
\noindent Hence we can view the collection of functions $\{K(\cdot, y) \}_{y \in E}$
as a subset of the product space
    %
    \[\{K(\cdot, y) \}_{y \in E} \subseteq \bigtimes_{x \in E} [0, C_x] \subseteq \R^E_{\geq0},\]
    %
\noindent so $\{K(\cdot, y) \}_{y \in E}$ is precompact under the product topology.
For the graded chains considered here, the Martin kernel separates states, so we can identify $E$ with
$\{K(\cdot, y) \}_{y \in E}$. The closure $\overline{E} \subseteq \R^E_{\geq 0}$ under this
identification is compact. We call $\overline{E}$ the \textit{Doob-Martin compactification}
of $E$ and its boundary $\partial E := \overline{E} \setminus E$ the \textit{Doob-Martin boundary}.

In summary, the map $y \mapsto K(\cdot,y)$ embeds $E$ into $\R^E_{\geq 0}$, and the
closure of its image in the product topology defines the Doob--Martin compactification.
Since $E$ is countable, this topology is metrizable. For example, after enumerating
$E = \{x_1,x_2,\ldots\}$, a compatible metric is
%
\[
d(y,z) := \sum_{i=1}^{\infty} 2^{-i}
\min\left\{1,\left|K(x_i,y)-K(x_i,z)\right|\right\}.
\]

%
\begin{remark}
    By construction, a sequence $\{y_n\} \in E$ converges to a point
    $\alpha \in \overline{E}$ if and only if $\{K(\cdot, y_n)\}$ converges
    in $\R^E_{\geq 0}$, namely if and only if $K(x, y_n)$ converges in $\R_{\geq 0}$
     for all $x \in E$. Thus for fixed $x$ the function
     $K(x, \cdot) : E \to \R_{\geq 0}$ extends continuously to $\overline{E}$.
    \label{rmk:conv_in_dm}
\end{remark}
%
%
\noindent There is a close relationship between the Doob-Martin kernel $K$ and
the backwards transition probabilities of the chain. Namely, for
$x \in E_n, y_N \in E_{N}$ Bayes rule implies
%
\[K(x,y_N) = \frac{\P(X_{N} = y_N \; | \; X_n = x)}{\P(X_{N} = y_N)} = \frac{\P(X_n = x \; | \; X_{N} = y_N)}{\P(X_n = x)}.\]
%
\noindent Thus $K(x, y_N)$ converges as $N \to \infty$ for all $x$ if and only
 if the conditional probabilities $\P(X_n = x \; | \; X_N = y_N)$ converge for
 all $x$.

%%%%%%%%%%%%%%%%%%%%%%%%%%%%
% Minimal Boundary and representation of harmonics
%%%%%%%%%%%%%%%%%%%%%%%%%%%%
\subsection{Minimal Boundary and Representation of Harmonics}

\noindent Given a Markov chain $X_n$ on state space $E$, we say a function $h:E \to \R$ is \textit{harmonic} if
%
    \[h(x) = \E [h(X_{n+1}) \; | \; X_n = x].\]
%
\noindent Note $K(\cdot, \alpha)$ is harmonic for every $\alpha \in \partial E$. The \textit{minimal boundary} is the subset
%
    \begin{equation}
        \partial_m E := \{ \alpha \in \partial E : K(\cdot, \alpha) \text{ is minimal harmonic}\}.
    \end{equation}
%
\noindent A \textit{minimal harmonic} function is a harmonic function $u \geq 0$ so
 for any other harmonic $w$ satisfying $0 \leq w \leq u$ we must have $w = c u$ for
 some constant $c \geq 0$.

The set $\Hn$ of nonnegative harmonics with $h(x_0) = 1$ is a compact, convex set whose
extreme points are precisely the minimal harmonics. Clearly $K(\cdot, \alpha) \in \Hn$, so
the Doob-Martin boundary can be thought of as a subset of $\Hn$, and
the minimal boundary is precisely the set of its extreme points.
Any nonnegative harmonic function $u$ can be represented
as an integral over these extreme points

\begin{equation}
    u(x) = \int_{\partial_m E} K(x, \alpha) \mu_u(d\alpha).
    \label{eqn:harmonic_representation_theorem}
\end{equation}

\noindent Moreover, as $K(\cdot, \alpha)$ is a harmonic for $\alpha \in \partial_m E$
and $|K(x, \alpha)| \leq C_x$, equation~\eqref{eqn:harmonic_representation_theorem}
defines a nonnegative harmonic function $u$ for any finite measure $\mu_u$, and $u$ satisfies
$u(x_0) = 1$ if and only if $\mu_u$ is a probability measure.

In words, nonnegative harmonic functions are characterized by their distribution on the minimal boundary $\partial_m E$.
For the special case of $u = K(\cdot, \alpha)$ for $\alpha \in \partial_m E$ clearly we have $\mu_{u} = \delta_\alpha$.
In some cases, such as the R\'emy chain studied in~\cite{evans2014_doob_martin_remy_tree_growth}, $\partial E = \partial_m E$, but in
general this is not always the case. A full proof of these facts can be found in~\cite{sawyer_martin_boundaries_random_walks}.

%%%%%%%%%%%%%%%%%%%%%%%%%%%%
% h processes
%%%%%%%%%%%%%%%%%%%%%%%%%%%%
\subsection{Using h-processes to Condition at Infinity} \label{sec:h_processes}

 Given a harmonic function $h \in \Hn$, we can define a new transition mechanism $p^h$
  which ``tilts'' the dynamics using $h$. For any $y \in E$ and $x \in E_h := \{x : h(x) > 0\}$ let
%
\begin{equation}
    p^h(x,y) := \frac{h(y)}{h(x)} p(x,y).
    \label{eqn:h-process}
\end{equation}
%
\noindent Thus the new chain reweights transitions of the original chain by a
factor of $h(y) / h(x)$, favoring transitions towards states where $h$ is large.
When $h \equiv 1$ we recover the original chain.

We call the new Markov chain $X^{(h)}_n$ with transition probabilities $p^h$ the
\textit{Doob h-process} or \textit{Doob h-transform} of $X_n$ corresponding to
$h(x)$; it has state space $E_h$. The $h$-processes $X_n^{(K(\cdot, \alpha))}$ are of natural interest
because they characterize the distributions of all other $h$-processes:
%
\begin{equation}
\P_{x_0}^h( B) = \int_{\partial_m E} \P_{x_0}^{K(\cdot, \alpha)}(B) \; \mu_h(d\alpha).
    \label{eqn:K_alpha_processes_char_h_processes}
\end{equation}%
%
\noindent for the unique measure $\mu_h$ supported on the minimal boundary as given in
\eqref{eqn:harmonic_representation_theorem}. In particular, these $h$-processes converge
almost surely in the Doob-Martin topology, and
%
\begin{equation}
    \P_{x}(\lim X_n^{(h)} \in A) = \frac{1}{h(x)} \int_{A \cap \partial_m E} K(x, \alpha) \mu_h(d\alpha),
    \label{eqn:limit_distribution_h-process}
\end{equation}

\noindent Since $K(x_0, \alpha) = h(x_0) = 1$, an $h$-process started at our basepoint $x_0$
is merely a chain conditioned to have distribution $\mu_h$ on the minimal boundary. Again, if $\alpha \in \partial_m E$
then $\mu_{K(\cdot, \alpha)} = \delta_\alpha$ so
%
    \[X_n^{K(\cdot, \alpha)} \xrightarrow{a.s.} \alpha.\]
%
Thus we can think of the chain $X^{(K(\cdot, \alpha))}_n$ to be our original Markov chain
conditioned to equal $\alpha \in \partial_m E$  ``at infinity''.


% Bridges and its relationship to the Doob martin boundary
\subsection{Infinite Bridges} \label{sec:infinite_bridges}

 \indent A finite-step Markov process $(Y_n)_{n = 1}^N$ with $Y_0 = x_0$ which
 has the same backwards transition probabilities as $(X_n)_{n \in \N}$ is called
 a (finite) \textit{bridge} for $X_n$
 %
    \[\P(X_n = i \mid X_{n+1} = j) = \P(Y_n = i \mid Y_{n+1} = j), \quad \forall i,j \in E, n < N.\]
 %
 \noindent Similarly, if the index set is $\N$ we call
  $(Y_n)_{n \in \N}$ an \textit{infinite bridge}. Let $(X^{(y)}_0, X^{(y)}_1, \ldots, X^{(y)}_{N(y)})$
  be the finite-step chain obtained from $X_n$ by conditioning on
  $X_{N(y)} = y$. Then
%
 \begin{align*}
     \P(X^{(y)}_n = i | X^{(y)}_{n+1} = j) & =  \frac{\P_{x_0}(X_n = i, X_{n+1} = j, X_{N(y)} = y)}{\P_{x_0}( X_{n+1} = j, X_{N(y)} = y)} \\
     %
     & =  \frac{\P_{x_0}(X_n = i, X_{n+1} = j) \cdot \P_{j}( X_{N(y)} = y)}{\P_{x_0}( X_{n+1} = j) \cdot \P_j( X_{N(y)} = y)} \\
     %
     & =  \P_{x_0}(X_n = i | X_{n+1} = j).
 \end{align*}
%
 Namely, $X^{(y)}_n$ is a finite bridge of $X_n$. Moreover,
%
 \begin{align*}
     \P(X^{(y)}_{n+1} = i | X^{(y)}_n = j) & =  \frac{\P_{x_0}(X_n = j, X_{n+1} = i, X_{N(y)} = y)}{\P_{x_0}( X_n=j, X_{N(y)} = y)} \\
     %
     & = \frac{\P_{x_0}(X \text{ hits j}) \cdot p(j,i) \cdot \P_i(X \text{ hits y})}{\P_{x_0}(X \text{ hits j}) \cdot \P_j(X \text{ hits y})} \\
     %
     & = K(j,y)^{-1}\cdot p(j,i) \cdot K( i, y).
 \end{align*}

\noindent This is essentially the definition of an $h$-process
with $h = K(\cdot, y)$, except for the fact that $K(\cdot, y)$ is not harmonic for $y \in E$.
However, as before, intuition suggests we ought to take the limit as $y$ ``goes to infinity'' to get a truly
harmonic function $K(\cdot, \alpha)$. Namely, by Remark~\ref{rmk:conv_in_dm} we know a
sequence $(y_k)$ in $E$ converges to a point $\alpha \in \partial E$ if and only if
 $K(x, y_k)$ converges in $\R_{\geq 0}$ to $K(x, \alpha)$ for all $x \in E$ which
 of course implies
 %
    \[\P(X^{(y_k)}_{n+1} = i | X^{(y_k)}_n = j) \xrightarrow{k\to \infty} K(j,\alpha)^{-1}\cdot p(j,i) \cdot K( i, \alpha) =: p^{K(\cdot, \alpha)}(j,i).\]
%

\noindent As observed in~\cite{follmer1975_phase_transition_martin_boundary}, such a sequence converges in the Doob
Martin topology if and only if finite initial segments of the corresponding bridges converge
 in distribution, namely
 %
 \begin{equation}
    (X^{(y_k)}_0, X^{(y_k)}_1, \ldots, X^{(y_k)}_M) \Rightarrow (X^{K(\cdot, \alpha)}_0, X^{K(\cdot, \alpha)}_1, \ldots, X^{K(\cdot, \alpha)}_M)
    \label{eqn:2.3_initial_sections_finite_bridges}
\end{equation}
 %
   as $k \to \infty$, for every finite $M$. More generally, given any nonnegative harmonic
   function with $h(x_0) = 1$, its Doob h-transform is an infinite bridge
   %
        \begin{align*}
        \P(X^h_n = x \; | \; X^h_{n+1} = y) & = \frac{\P(X^h_{n+1} = y \; | \; X^h_{n} = x) \cdot \P(X^h_n = x)}{\P(X^h_{n+1} = y)} \\
        %
        & = \frac{\frac{h(y)}{h(x)}\P(X_{n+1} = y \; | \; X_{n} = x) \cdot \frac{h(x)}{h(x_0)}\P(X_n = x)}{\frac{h(y)}{h(x_0)}\P(X_{n+1} = y)} \\
        %
        & = \P(X_n = x \; | \; X_{n+1} = y).
        \end{align*}
    %


\noindent Conversely, given any infinite bridge $X^\infty_n$, then
\eqref{eqn:K_alpha_processes_char_h_processes} tells us there
 exists a measure $\mu_h$ so that $X^\infty_n$ is a mixture of the extremal
  bridges $\P^{K(\cdot, \alpha)}$ and the corresponding harmonic $h$ is given by
  \eqref{eqn:harmonic_representation_theorem}.
%

Call an infinite bridge \textit{extremal} if it has a trivial tail
$\sigma$-field, or equivalently if it cannot be written as a nondegenerate
mixture of other infinite bridges. By~\eqref{eqn:K_alpha_processes_char_h_processes}
and the correspondence between bridges and
harmonic functions, all infinite bridges are just mixtures of the extremal
bridges. Hence characterizing the extremal bridges allows us to
characterize all infinite bridges.

Our observations from Section~\ref{sec:h_processes} tell us that for $\alpha \in \partial_m E$
the conditioned chain $X_n^{K(\cdot, \alpha)}$ is extremal as it converges almost
 surely to the constant $\alpha$. Moreover, the converse is true: $X^{K(\cdot, \alpha)}$
 is extremal only if $\alpha \in \partial_m E$. Hence characterizing the extremal
 bridges will allow us to characterize the minimal boundary, and thus by
 \eqref{eqn:harmonic_representation_theorem} characterize all nonnegative
  harmonic functions for this chain.

In summary, points of the Doob--Martin boundary can be described as limits of
sequences in $E$, or equivalently through limits of the corresponding finite
bridges as in~\eqref{eqn:2.3_initial_sections_finite_bridges}. Within this
boundary, the points of the minimal boundary correspond to extremal infinite
bridges and extremal nonnegative harmonic functions.



% Proof Appendix %%%%%%%%%%%%%%%%%%%%%%%%%%%%%%%%%%%
\newpage
\section{Proofs}
\label{appendix:proof_appendix}

\subsection{Proofs of Section~\ref{sec:exchangeable_array_processes} (Exchangeable Array Processes)}

\noindent\textbf{Lemma~\ref{lem:state_equivalence} (State-Space Projections of Exchangeable Arrays).}
If $(X_e)$ is an exchangeable array with entries in $S$ and $g:S \to S'$ is
measurable, then $(g(X_e))$ is an exchangeable array of the same dimension.
Moreover, if $f$ is the representing function of $(X_e)$, then $g \circ f$ is
the representing function of $(g(X_e))$.

\begin{proof}[Proof of Lemma~\ref{lem:state_equivalence} (State-Space Projections of Exchangeable Arrays)]
The result follows immediately from applying $g$ coordinatewise. If $(X_e)$ is exchangeable, then for each finite permutation $\sigma$,
\[
(g(X_{\sigma(e)}))_{e \in \binom{\N}{d}} \overset{\mathrm{law}}{=} (g(X_e))_{e \in \binom{\N}{d}}.
\]
Moreover, if $X_e = f(U_\emptyset,(U_i)_{i \in e},\ldots,U_e)$ in distribution, then $g(X_e)$ is represented by $g \circ f$.
\end{proof}

\noindent\textbf{Lemma~\ref{lem:index_equivalence} (Markov Property of the Unlabeled Chain).}
When the state space $S$ is countable, the process $M_n$ is a Markov chain.

\begin{proof}[Proof of Lemma~\ref{lem:index_equivalence} (Markov Property of the Unlabeled Chain)]
First, note that the definition of an exchangeable array implies
%
\[\P(L_n = x) = \P(L_n = x') \qquad \forall x \sim x'.\]

\noindent Moreover if $x \sim x'$ are related by $\sigma \in \operatorname{Sym}(\N)$ then again exchangeability implies
%
\begin{align*}
    \P(M_{n+1} = [y]_\sim,\; L_n = x) & = \sum_{y' \sim y} \P(L_{n+1} = y', \; L_n = x)\\
    & = \sum_{y' \sim y} \P(L_{n+1} = \sigma(y'), \; L_n = \sigma(x)) \\
    & = \sum_{z \sim y} \P(L_{n+1} = z, \; L_n = x') \\
    & = \P(M_{n+1} = [y]_\sim,\; L_n = x').
\end{align*}
%
\noindent Thus we see
%
\[\P(M_{n+1} = [y]_\sim \; | \; L_n = x) = \P(M_{n+1} = [y]_\sim \; | \; L_n = x') \qquad \forall x \sim x', \quad [y]_\sim \in [E_{n+1}]_\sim.\]
%
\noindent Thus the transition probability from $x$ into the class $[y]_\sim$
depends only on the equivalence class $[x]_\sim$. The projected transition
kernel is therefore well-defined on equivalence classes, and hence $M_n$ is a
Markov chain.
\end{proof}

\subsection{Proofs of Section~\ref{sec:exchangeable_array_boundary} (Doob-Martin Boundary)}

\noindent\textbf{Lemma~\ref{lem:boundary_labeled_chain} (Boundary of the Labeled Chain).}
The boundary of $L_n$ consists of the deterministic arrays
$x \in \mathrm{supp}(X_e)$.

\begin{proof}[Proof of Lemma~\ref{lem:boundary_labeled_chain} (Boundary of the Labeled Chain)]\leavevmode\par\noindent
Fix $x\in E_m$ and $y\in E_n$ with $n\geq m$. Since the backwards dynamics
of the labeled chain are deterministic,
\[
\P(L_m=x\mid L_n=y)=\ind\{y|_{[m]}=x\}.
\]
Consequently, the Doob--Martin kernel of the labeled chain is
\[
K_L(x,y)=\frac{\ind\{y|_{[m]}=x\}}{\P(L_m=x)}.
\]
Thus a sequence $(y_n)$ with $y_n\in E_n$ converges in the Doob--Martin
topology if and only if, for every $m$, the restriction $y_n|_{[m]}$ is
eventually constant. These limiting restrictions are consistent and therefore
define a unique deterministic infinite array in $\operatorname{supp}(X_e)$.
Conversely, the finite restrictions of any deterministic array in
$\operatorname{supp}(X_e)$ form a convergent sequence with that boundary
point as its limit.
\end{proof}

\noindent\textbf{Lemma~\ref{lem:harmonic_functions_labeled} (Exchangeable Bridges of the Labeled Chain).}
Any infinite bridge of the labeled exchangeable array process is an array
process, and it is exchangeable if and only if the corresponding harmonic
function $h$ is constant on the equivalence classes of $\sim$.

\begin{proof}[Proof of Lemma~\ref{lem:harmonic_functions_labeled} (Exchangeable Bridges of the Labeled Chain)]
Suppose the harmonic $h \geq 0$ with $h(x_0) = 1$ is the one corresponding to this particular bridge. If $L_n$ is the original labeled exchangeable array process and $L^h_n$ is the infinite bridge
%
\[\P(L^h_0 = x_0, \ldots, L^h_n = x_n) = h(x_n) \cdot \P(L_0 = x_0, \ldots, L_n = x_n).\]
%
\noindent The right hand side is zero if the states $x_0, \ldots , x_n$ are not nested in the sense of arrays, $x_m = (x_n) |_{\binom{[m]}{d}}$ for $m < n$, and thus the same is true of the left hand side. Thus, by taking a projective limit, the chain $L_n^h$ a.s. defines an array ($X^h_e)$ of the same dimension as the original exchangeable array $(X_e)$ of which it is the corresponding array process.

Now if $\sigma \in \operatorname{Sym}(\N)$ let $\tx := \sigma(x)$ be the array generated by permuting the coordinates of $x$. Then the definition of the Doob h-transform and the exchangeability of the original array imply
%
\begin{align*}
    \P(L^h_0 = \tx_0, \ldots, L^h_n = \tx_n)  & = h(\tx_n) \cdot \P(L_0 = \tx_0, \ldots, L_n = \tx_n) \\
    %
    & = h(\tx_n) \cdot \P(L_0 = x_0, \ldots, L_n = x_n) \\
    & = \frac{h(\tx_n)}{h(x_n)} \cdot \P(L^h_0 = x_0, \ldots, L^h_n = x_n).
\end{align*}

\noindent These two are equal if and only if $h(\tx_n) = h(x_n)$. Since $x_n$ and $\sigma$ were arbitrary, the result follows from the fact that the finite dimensional distributions of our array $(X^h_e)$ are determined from the distribution of $(L_n^h)$.
\end{proof}

\noindent\textbf{Corollary~\ref{cor:unlabeled_associated_to_exchangeable_array} (Bridges of the Unlabeled Chain).}
If $M^{\tilde{h}}_n$ is an infinite bridge of the unlabeled chain $M_n$, then
$M^{\tilde{h}}_n$ equals the unlabeled process associated to some exchangeable
array $(X_e^h)$. In particular, if $\tilde{h}$ is the harmonic corresponding to
the bridge $M^{\tilde{h}}_n$, then $h(x) := \tilde{h}([x])$ is the harmonic
corresponding to the array $(X_e^h)$.

\begin{proof}[Proof of Corollary~\ref{cor:unlabeled_associated_to_exchangeable_array} (Bridges of the Unlabeled Chain)]
The bridges of $(M_n)$ are in correspondence with harmonic functions
$\tilde{h}$ of this chain. Define $h(x) := \tilde{h}([x])$ on the labeled state
space. Since the transition kernel of $M_n$ is obtained by summing the
transition probabilities of $L_n$ over equivalence classes, harmonicity of
$\tilde{h}$ gives
%
\[
    h(x)
    =
    \tilde{h}([x])
    =
    \sum_{[y]_\sim} p_M([x]_\sim,[y]_\sim)\tilde{h}([y]_\sim)
    =
    \sum_y p_L(x,y) h(y).
\]
%
Thus $h$ is a harmonic function for the labeled chain and is constant on the
equivalence classes of $\sim$. Applying Lemma~\ref{lem:harmonic_functions_labeled}
gives the corollary.
\end{proof}

\noindent\textbf{Theorem~\ref{thm:boundary_is_ergodic_arrays} (Extremal Bridges and Ergodic Arrays).}
For any infinite bridge of the unlabeled exchangeable array process $M_n$, the
following are equivalent:
\begin{enumerate}
    \item The infinite bridge is extremal.
    \item Its corresponding exchangeable array is ergodic.
    \item Its corresponding exchangeable array is dissociated.
    \item Its corresponding exchangeable array admits a representation function
    $f$ that does not depend on $U_\emptyset$ almost surely.
\end{enumerate}

\begin{proof}[Proof of Theorem~\ref{thm:boundary_is_ergodic_arrays} (Extremal Bridges and Ergodic Arrays)]
The last three statements are equivalent by Lemma 7.35 of~\cite{kallenberg2005_probabilistic_symmetries}. For the implication $(2) \to (1)$, suppose for contraposition that we have a nonextremal bridge $M_n^{\tilde{h}}$ associated to the harmonic $\tilde{h}$. Then equation~\eqref{eqn:K_alpha_processes_char_h_processes} applied to the Doob-Martin kernel $\tilde{K}$ of our unlabeled chain implies

\[\P((M_n^{\tilde{h}}) = [x]_\sim) = \int_{\partial_m \tilde{E}} \P^{\tilde{K}(\cdot, \alpha)}([x]_\sim) \; \mu_{\tilde{h}}(d\alpha),\]

\noindent for some nondegenerate measure $\mu_{\tilde{h}}$ on the minimal boundary $\partial_m \tilde{E}$ of the unlabeled chain. By construction of the Doob-Martin boundary, the functions $\tilde{K}(\cdot, \alpha)$ are harmonic for the unlabeled chain. Thus by Corollary~\ref{cor:unlabeled_associated_to_exchangeable_array} the distributions $\P^{\tilde{K}(\cdot, \alpha)}$ correspond to the unlabeled chains associated to exchangeable arrays $(X^\alpha_e)$. Likewise, $M_n^{\tilde{h}}$ corresponds to some exchangeable array $(X_e^{\tilde{h}})$ as well. Thus we can rephrase the previous equation in terms of the associated exchangeable arrays

\[\P((X^{\tilde{h}}_e) \in [x]_\sim) = \int_{\partial_m \tilde{E}} \P(X_e^\alpha \in [x]_\sim) \; \mu_{\tilde{h}}(d\alpha).\]

\noindent These probabilities are constant on the equivalence classes by exchangeability. Hence

\[\P((X^{\tilde{h}}_e) = x) = \int_{\partial_m \tilde{E}} \P(X_e^\alpha = x) \; \mu_{\tilde{h}}(d\alpha).\]

\noindent Thus as $\mu_{\tilde{h}}$ is nondegenerate, $(X_e^{\tilde{h}})$ is a nondegenerate mixture of other exchangeable arrays, and hence not ergodic.

For the implication $(1) \to (2)$, suppose for contraposition that the
exchangeable array associated with a bridge $M_n^\alpha$ has a nonergodic law
$\nu$. Then there is a $\nu$-invariant event $I$ with $0<\nu(I)<1$. Define
the conditional laws
\[
    \nu_1(\,\cdot\,) := \nu(\,\cdot\mid I),
    \qquad
    \nu_0(\,\cdot\,) := \nu(\,\cdot\mid I^c).
\]
Because $I$ is invariant and $\nu$ is exchangeable, both $\nu_0$ and $\nu_1$
are exchangeable. Let $M_n^{(0)}$ and $M_n^{(1)}$ be their associated
unlabeled processes.

We claim that both processes are infinite bridges of $M_n$. Indeed, under any
exchangeable array law, conditioning on $M_{n+1}=[y]_\sim$ makes the labeled
array at level $n+1$ uniform over $[y]_\sim$. Consequently,
\[
    \P_{\nu_i}\left(M_n=[x]_\sim\mid M_{n+1}=[y]_\sim\right)
    =
    \frac{\left|\left\{z\in[y]_\sim:z|_{[n]}\in[x]_\sim\right\}\right|}
         {|[y]_\sim|},
    \qquad i\in\{0,1\}.
\]
This backwards transition probability depends only on the equivalence classes
$[x]_\sim$ and $[y]_\sim$, not on the particular exchangeable law. Hence
$M_n^{(0)}$ and $M_n^{(1)}$ have the same backwards transitions as
$M_n^\alpha$, and therefore as the original chain $M_n$.

Finally, projection from a labeled array to its unlabeled process is affine, so
\[
    \mathcal{L}\bigl((M_n^\alpha)_{n\geq0}\bigr)
    = \nu(I)\,\mathcal{L}\bigl((M_n^{(1)})_{n\geq0}\bigr)
      +(1-\nu(I))\,\mathcal{L}\bigl((M_n^{(0)})_{n\geq0}\bigr).
\]
This mixture is nontrivial. Indeed, if the two unlabeled process laws were
equal, exchangeability and uniformity within each equivalence class would
imply that all their finite labeled distributions were equal, and hence that
$\nu_0=\nu_1$. This is impossible because $\nu_1(I)=1$ whereas
$\nu_0(I)=0$. Thus $M_n^\alpha$ is not extremal, proving the contraposition.
\end{proof}

\subsection{Proofs of Section~\ref{sec:embedding_density_representation_functions} (Embedding Densities)}

\noindent\textbf{Lemma~\ref{lem:relating_different_notions_of_embedding} (Exact and Weak Embedding Densities).}
For binary arrays $x,y$, the following hold:
\begin{enumerate}[label=(\roman*)]
    \item $emb(x,y) \leq emb^W(x,y)$.
    \item $emb^W(x,y) = \sum_{x' \geq x} emb(x',y)$.
    \item $emb(x,y) = \sum_{x' \geq x} (-1)^{|\ind(x')|-|\ind(x)|} \cdot emb^W(x',y)$.
\end{enumerate}

\begin{proof}
These proofs extend analogous arguments from graphon theory in~\cite{meliot2024_graphons_permutons_survey}.
First, (i) is immediate because every embedding is a weak embedding. For (ii), note that if
$\sigma \in \operatorname{Sym}(\N)$ is such that $x \lesssim \sigma(y)$,
 then there is precisely one $x' \geq x$ so that $x' \leq y$, namely $x'$ equals the array
 $\sigma(y)$ restricted to $\binom{\{1,\ldots,|x|\}}{d}$. For (iii), if for each $e$ where $x_e=0$ we let $A_e$ be the set of weak embeddings of $x'=x+\delta_e$ into $y$, then
%
\[\{\text{embeddings of $x$ into $y$}\} = \{\text{weak embeddings of $x$ into $y$}\} \setminus \bigcup_{e \;:\; x_e=0} A_e.\]
%
\noindent The formula follows from inclusion--exclusion.
\end{proof}

\noindent\textbf{Lemma~\ref{lem:equivalence_of_dm_and_density_convergence} (Doob-Martin Convergence and Embedding Densities).}
Let $y_N\in E_N$. The following are equivalent:
\begin{enumerate}[label=(\roman*)]
    \item The sequence $y_N$ converges in the Doob-Martin topology.
    \item The embedding density $emb(x, y_N)$ converges for all $x$.
\end{enumerate}
For binary arrays, convergence of the weak embedding density is also
equivalent.

\begin{proof}
Exchangeability implies

\[\P(L_n = x \; |\; M_n = [x]_\sim) = \frac{\P(L_n = x)}{\P(M_n = [x]_\sim) }=\frac{1}{| \; [x]_\sim \; |}.\]

\noindent As $L_n$ has deterministic backwards dynamics, this implies the unlabeled chain $M_n$
has
%
    \[\P(M_{n} = [x]_\sim \mid M_{n+m} = [y]_\sim) = \frac{| \{z \in [y]_\sim : z|_{[n]} \in [x]_\sim\} \; |}{|\; [y]_\sim \; |} \quad x \in E_n, y \in E_{n+m},\]
%
\noindent where $z|_{[n]}$ is the restriction of $z$ to indices in $\binom{[n]}{d}$. Since any
permutation of $[n]$ extends to a permutation of $[n+m]$, the number of elements
$z\in[y]_\sim$ satisfying $z|_{[n]}=\widetilde x$ is the same for every $\widetilde x\in[x]_\sim$.
Thus the above equals
%
    \[= \frac{|\{z \in [y]_\sim : z\mid_{[n]} = x \}| }{|[y]_\sim|} \cdot |[x]_\sim| = emb(x,y) \cdot |[x]_\sim|.\]
%
\noindent As a result, we have Doob-Martin kernel for the unlabeled chain
%
    \[K([x]_\sim, [y]_\sim) = \frac{emb(x, y) |[x]_\sim|}{\P(M_n = [x]_\sim)} = \frac{emb(x,y)}{\P(L_n = x)}.\]
%
\noindent Hence the embedding density and Doob-Martin kernel differ only up to a constant depending on $x$.
Thus for fixed $x$, one converges if and only if the other does.
For binary arrays, Lemma~\ref{lem:relating_different_notions_of_embedding}
shows that convergence of all exact embedding densities is equivalent to
convergence of all weak embedding densities.
\end{proof}

\noindent\textbf{Lemma~\ref{lem:integral_representations_of_embedding_densities} (Embedding Densities of Representation Functions).}
Suppose $d \leq k \leq n$. If $F$ is a size-$k$, $d$-dimensional array and
$f$ is the representing function of an exchangeable array $X$, then
\begin{equation*}
    emb(F,f) = \int_{[0,1]^{2^k}}  \prod_{S \in \binom{[k]}{d}} \ind \{f(x_{(S)}) = F_S\} \; d\lambda_{2^k}.
\end{equation*}
If $F$ and $f$ are binary-valued, then also
\begin{equation*}
    emb^W(F,f) = \int_{[0,1]^{2^k}}  \prod_{S \in \ind(F)} f(x_{(S)}) \; d\lambda_{2^k}.
\end{equation*}
These quantities are independent of $n$. Moreover, if $F$ and $f$ are binary-valued and $f$ does not depend on $x_\emptyset$,
\begin{align*}
    \P\left( |emb^W(F, X_n(f)) - emb^W(F, f)| > \epsilon \right)
    &\leq 2 \exp \left(-\frac{\epsilon^2}{2k^2} n \right),\\[5pt]
    \P\left( |emb(F, X_n(f)) - emb(F, f)| > \epsilon \right)
    &\leq 2 \exp \left(-\frac{\epsilon^2}{2k^2} n \right).
\end{align*}
Thus
\[
\lim_{n \to \infty} emb(F, X_n(f)) \overset{a.s.}{=} emb(F,f),
\]
and
\[
    \lim_{n \to \infty} emb^W(F, X_n(f)) \overset{a.s.}{=} emb^W(F,f).
\]
In particular, Lemma~\ref{lem:relating_different_notions_of_embedding} (i)--(iii)
extends to $emb^W(F,f)$ and $emb(F,f)$.

\begin{proof}[Proof of Lemma~\ref{lem:integral_representations_of_embedding_densities} (Embedding Densities of Representation Functions)]\leavevmode\par\noindent
First we prove the expectation formulas, starting with $emb^W$. Let
$\phi : [k] \to [n]$ be injective. Then
    \begin{align*}
        \P(\text{$\phi$ is an array embedding}) &= \int_{[0,1]^{2^n}} \prod_{S \in \ind(F)} f(x_{(\phi(S))}) \, d\lambda_{2^n}
        \intertext{By the injectivity of $\phi$ we can apply a change of coordinates,}
        &= \int_{[0,1]^{2^n}} \prod_{S \in \ind(F)} f(x_{(S)}) \, d\lambda_{2^n} \\
        \intertext{After integrating out all the superfluous coordinates,}
        &= emb^W(F,f).
    \end{align*}
    So, writing $emb^W(F,X_n(f))$ in terms of the indicator functions $\ind_{\{\text{$\varphi$ is an array embedding}\}}$,
    \begin{equation*}
        \E[emb^W(F,X_n(f))] = emb^W(F,f).
    \end{equation*}

    \noindent The proof for $emb$ is identical.

    For the concentration results, we follow exactly the proof of Theorem 2.5 in~\cite{lovasz2004_limits_dense_graph_sequences},
    except that we choose a different filtration meant to capture the randomness available at different levels of a
    multidimensional array. Namely, define filtration
    %
        \[\cF_m := \sigma\left(\left\{U_A : A \subseteq [n], \max A \leq m\right\}\right), \quad \cF_0 := \sigma(U_\emptyset).\]
    %
    \noindent This can be used to define a Doob martingale
    %
        \[B_m := \E (emb^W(F, X_n(f)) \mid \cF_m).\]
    %
    \noindent Note that
    %
        \[B_0 = \E (emb^W(F, X_n(f)) \mid U_\emptyset) = emb^W(F, f_{U_\emptyset}), \quad B_n = emb^W(F, X_n(f)).\]
    %
    \noindent Here $f_{U_\emptyset}$ is the (random) representation function
    %
        \[f_{U_\emptyset}(u) := f(U_\emptyset, (u_i)_{i \in [d]}, \ldots, u_{[d]}).\]
    %
    \noindent If $f$ does not depend on $U_\emptyset$, then $f_{U_\emptyset}=f$ is deterministic.
    Bounding the increments using a bound on the number of injective functions $\phi:[k] \to [n]$
    and applying Azuma's inequality as in~\cite{lovasz2004_limits_dense_graph_sequences} gives the desired inequalities.
    Then apply an identical procedure for $emb$.

    The convergence then follows immediately from Borel-Cantelli.
\end{proof}

\noindent\textbf{Lemma~\ref{lem:sampling_functions_approximate_embedding_densities}
(Sampling Functions Approximate Weak Embedding Densities).}
Let $F$ and $z$ be finite binary $d$-dimensional arrays with $|F|\leq|z|$.
Then
\[
    \left|emb^W(F,f_z)-emb^W(F,z)\right|
    \leq \frac{1}{|z|}\binom{|F|}{2}.
\]

\begin{proof}
Let $\Phi:[|F|]\to[|z|]$ be a uniformly random map. By the definition of
$f_z$ and the integral formula for weak embedding densities,
$emb^W(F,f_z)$ is the probability that
\[
    \{\Phi(i):i\in S\}\in\ind(z)
    \qquad\text{for every }S\in\ind(F).
\]
Conditional on $\Phi$ being injective, it is a uniformly random injection, so
the conditional probability of this event is $emb^W(F,z)$. The difference
between the unconditional and conditional probabilities is therefore at most
\[
    \P(\Phi\text{ is not injective})
    \leq \frac{1}{|z|}\binom{|F|}{2},
\]
where the last inequality is the union bound over pairs of indices in $[|F|]$.
\end{proof}

\end{subappendices}
